%% Introduction du livre — écrite par l'agent Architecte.
%% N'utilise QUE des commandes sémantiques bookforge.
\introLivre{
  \paragraphe{Tu sais coder. Tu lis un diff sans broncher, tu connais ton
  dépôt, tu as tes habitudes. Mais un \terme{agent} de développement, tu n'en
  as jamais vraiment lancé un dans ton propre projet — et l'idée de lui laisser
  modifier ton code te met mal à l'aise. C'est exactement pour toi que ce livre
  est écrit. Tu n'es pas débutant en programmation : tu es débutant en agents.}

  \paragraphe{La promesse est simple. À la fin, tu seras passé de « je n'ai
  jamais osé déléguer une tâche » à « j'orchestre \terme{Claude Code} dans mon
  travail quotidien avec confiance ». Pas une confiance aveugle : une confiance
  \important{outillée}, fondée sur ta capacité à déléguer la bonne tâche, à
  vérifier le résultat, et à poser des garde-fous.}

  \paragraphe{Ce livre ne reformule pas la documentation. On y délègue des
  tâches réelles, on relit ce que l'agent produit, on le prend en flagrant
  délit d'erreur, on surveille la facture et on désamorce les pièges avant
  qu'ils ne coûtent cher. L'objectif assumé est de battre les guides
  génériques qui se contentent de te dire de « mieux prompter ».}

  \paragraphe{La structure suit ta progression. Les premiers chapitres
  t'installent et te font déléguer ta toute première tâche. Le cœur du livre
  t'apprend à penser en tâches délégables, à piloter le \terme{contexte}, à
  vérifier le travail et à poser des garde-fous. Les derniers chapitres ouvrent
  sur l'automatisation, les outils connectés et l'orchestration de plusieurs
  agents, avant d'ancrer tout cela dans une routine quotidienne réaliste.}

  \paragraphe{Un mot sur le temps qui passe. Un outil de 2026 évolue vite : les
  commandes exactes, les versions et les tarifs d'aujourd'hui seront périmés
  demain. On a donc \important{isolé ce qui change vite} dans des chapitres
  clairement datés et faciles à réviser, et investi le reste du livre dans les
  principes durables : comment penser, déléguer et vérifier avec un agent. Ces
  réflexes-là te resteront, quelle que soit la version de l'outil.}
}
